% Chapter questions for sections/05_otp_transposition.tex, 05b_classical_extra.tex, 06_stego_math.tex.
% Numbers checked by script (scratchpad/r3/). The long practice set for these chapters is in questions/99_practice.tex.
\section{One-Time Pad, Transposition, Remaining Classical Topics and Steganography}\label{q:05_otp}
\noindent{\small Notes: \S\ref{sec:otp}, \S\ref{sec:transp}, \S\ref{sec:classextra}, \S\ref{sec:stego}, \S\ref{sec:primesintro}. Longer numericals on these chapters are in the practice set (Part III, sets F and G).}

\subsection*{Check your understanding}
\begin{enumerate}[leftmargin=1.6em,itemsep=3pt]
  \item \textbf{Define perfect secrecy and show that the one-time pad has it.} A cipher has perfect secrecy if $\Pr[M \mid C] = \Pr[M]$ for every message $M$ and ciphertext $C$, i.e.\ the ciphertext reveals nothing about the plaintext. In an OTP exactly one key, $K = C - M$, maps any $M$ to a given $C$, and all $26^n$ keys are equally likely, so $\Pr[C \mid M] = 1/26^n$ for every $M$; hence $\Pr[C] = 1/26^n$ and Bayes' rule gives $\Pr[M \mid C] = \Pr[M]$.
  \item \textbf{Why can no cipher with a key shorter than the message have perfect secrecy?} With fewer keys than messages, decrypting a given $C$ under every key yields only some of the possible messages, so the others get $\Pr[M \mid C] = 0 \ne \Pr[M]$; the ciphertext has ruled them out.
  \item \textbf{State the two fundamental difficulties of the one-time pad. Where is it still used?} (1) The practical problem of making large quantities of random keys; (2) the even more daunting problem of key distribution and protection. Because of them it is used only on low-bandwidth channels that need very high security, with the key carried in advance by a trusted courier.
  \item \textbf{What happens if a one-time-pad key is used for two messages?} The key cancels: $C_1 - C_2 = P_1 - P_2$ (or $C_1 \oplus C_2 = P_1 \oplus P_2$), which is two English texts mixed together and can be attacked; if one plaintext is known or guessed, $K = C_1 - P_1$ and the other message follows at once.
  \item \textbf{Why does a transposition cipher keep the letter frequencies, and how is this used?} It only moves letters, so the ciphertext is an anagram of the plaintext with the same letter counts. A ciphertext whose frequency count looks like ordinary English (many \ttx{e}, \ttx{t}, \ttx{a}) is therefore a transposition, not a substitution.
  \item \textbf{A 23-letter message is enciphered with a rail fence of depth 3. How long is each row?} The zig-zag repeats every $2(D-1) = 4$ letters and $23 = 5 \cdot 4 + 3$; each cycle gives row 1 and row 3 one letter and row 2 two, and the 3 leftover positions are rows 1, 2, 3. Rows: \textbf{6, 11, 6}.
  \item \textbf{Why is a double (two-stage) transposition stronger than a single columnar one?} After one stage, adjacent ciphertext letters come from plaintext positions a constant distance apart (the number of columns), so guessing the row count and rearranging columns breaks it (pattern attack). After two stages no constant step remains, so the pattern attack fails; the letter frequencies are still preserved, since the result is still one permutation.
  \item \textbf{Repeated trigrams in a Vigen\`ere ciphertext lie 35, 63 and 42 letters apart. What is the key length, and why?} Equal plaintext fragments encrypt identically only when their distance is a multiple of the key length $m$, so $m \mid \gcd(35, 63, 42) = 7$. Since 7 is prime and $m = 1$ would be a plain Caesar cipher, $m = 7$.
  \item \textbf{How does the index of coincidence confirm a key length?} The IC is the probability that two random letters of a text are equal: about 0.066 for English (also after any monoalphabetic substitution) and 0.038 for random text. Splitting a Vigen\`ere ciphertext into $m$ pieces gives pieces with IC near 0.066 only for the right $m$ (or its multiples), because only then is each piece a single Caesar cipher.
  \item \textbf{Why does the affine cipher have 312 keys, and why is $k_1 = 13$ not allowed?} $k_1$ must be invertible mod 26, i.e.\ in $\mathbb{Z}_{26}^*$ (12 values), and $k_2$ can be any of 26 values; different pairs give different substitutions, so $12 \times 26 = 312$. Since $\gcd(13, 26) = 13$, multiplication by 13 sends every even letter value to 0 and every odd one to 13, so it cannot be inverted.
  \item \textbf{An affine cipher maps \ttx{h} to \ttx{A} and \ttx{i} to \ttx{H}. Find the key.} $7k_1 + k_2 \equiv 0$ and $8k_1 + k_2 \equiv 7$; subtracting, $k_1 \equiv 7$ (valid, $\gcd(7,26) = 1$), and $k_2 \equiv -49 \equiv 3$. Key $\mathbf{(7, 3)}$; check: $8 \cdot 7 + 3 = 59 \equiv 7 = $ \ttx{H}.
  \item \textbf{Why does the Hill cipher resist frequency analysis but fall to a known-plaintext attack?} Each ciphertext letter depends on the whole $m$-letter block, so single-letter, digram and trigram statistics are hidden, and the key space $26^{m^2}$ is large. But the cipher is linear, $C = PK$, so $m$ known block pairs stacked into an invertible matrix $P$ give $K = P^{-1}C$; if $P$ is not invertible mod 26, other pairs are used.
  \item \textbf{A keyed transposition has encryption key (5 3 1 4 2). Find the decryption key.} DecKey[EncKey[$i$]] $= i$: DecKey[5] $= 1$, DecKey[3] $= 2$, DecKey[1] $= 3$, DecKey[4] $= 4$, DecKey[2] $= 5$, so the decryption key is \textbf{(3 5 2 4 1)}.
  \item \textbf{Why does Enigma never encrypt a letter to itself, and why did that matter?} The reflector sends the current back through the rotors on a different contact, so the output is always a different letter. Cryptanalysts used this to position guessed plaintext words (cribs) against the ciphertext, discarding every alignment in which a letter would map to itself.
  \item \textbf{Differentiate steganography from cryptography, give its two drawbacks, and say how to combine the two.} Cryptography hides the meaning of a message, steganography hides its existence. Steganography needs a lot of cover for few hidden bits, and once the method is discovered it is virtually worthless. Encrypting first and then hiding gives both: the message is unnoticed, and if found it is still ciphertext.
  \item \textbf{How many bytes can a $640 \times 480$ 24-bit image hide using one LSB per colour channel, and how much does a pixel change?} $640 \times 480 \times 3 = 921{,}600$ bits $= 115{,}200$ bytes; each channel value changes by at most 1, which is invisible.
\end{enumerate}
